% ============================================================
% v2/sections/intro.tex  はじめに（短縮・再構成版 2026-10-27）
% 方針：論文_再構成案.md（改訂版）。主の貢献は「成立条件の切り分け」。A は補助1（可能性）。
% ============================================================
\section{はじめに}\label{sec:intro}

シミュレーション説は，他者の立場に自分を置き，自分の心の働きを用いて他者の心的状態や行動を推し量ると考える\cite{gordon1986folk,goldman2006simulating}．
予測符号化の観点からは，自分の意思決定や計画の仕組みが，他者の行動を説明する仮説を試すためにも使われ，そこでの信念や欲求は自他を区別しない形で表されると論じられている\cite{gordon2021simulation}．
しかし，自分の価値に基づく予測の仕組みが，どのような条件で他者の予測に使われるようになるのかを，学習の過程から具体的に示した計算モデルは少ない．

Noguchi ら\cite{noguchi2022superposition}の重ね合わせ機構（superposition mechanism）は，一つの感覚入力を二つの経路で符号化し，重みを共有した一つのモジュールで処理する予測学習の仕組みである．
その新しさは，他者の動きを自分のどの経験に当てはめるかという自他の対応を与えずに，自分の次の視界を予測する学習だけから，自他の区別，他者の位置，他者の視点，他者の運動が生まれた点にある．
一方で，この機構が扱ったのは他者の位置・視点・運動までであり，他者の欲求のような内面には及んでいない\cite{noguchi2022superposition}．

本研究は，重ね合わせ機構の上で，自分の価値に従って動いた経験から学ばれた予測が他者の予測に使われるかを調べる．
ここで価値は，方向ごとの快・不快として人が与える信号である（\ref{sec:env}節）．
問うのは価値そのものが生まれるかではなく，与えた価値の信号とそれに従う行動の規則性を，自他で共有された予測学習がどう利用するかである．
また本研究は，自他の対応を担う視覚エンコーダを，文献\cite{noguchi2022superposition}の設定で学習した重みに固定して与える．
自他を同じ符号で表す性質は本研究の新しさではなく，前提である．

重みを共有している以上，自分側で学んだことが他者側にも働くのは当然だ，と考えられるかもしれない．
しかし，他者側に届くのが「入力された快・不快の信号」そのものなのか，「その信号が自分の動きの何を予測するか」として学ばれた解釈なのか，その解釈が仕組みのどこに宿るのかは自明でない．
そこで本研究は，次の問いを立てる．
\begin{quote}
  自分の価値の経験から学んだ予測は，自他で共有された予測の仕組みを通して，どのような条件のもとで他者の予測に使われるか．
\end{quote}
ここで「使われる」（転用）とは，他者側に与えた価値の入力に応じて他者の予測が変わることを指し，他者の行動の予測が正確になることとは区別する．
これを，価値に従って動いた経験の量を変えた学習と，三つの対照で調べた．
第一に，自他で重みを共有しない対照（非共有条件）．
第二に，自分の動きの軌跡はまったく同じまま，快・不快の符号だけをエピソードごとに無作為にし，価値の信号と行動の対応だけを壊した対照（符号無作為条件）．
第三に，学習を終えた非共有のモデルに，自分側の重みを写す比較である．

本研究の貢献は次のとおりである．
\begin{description}
  \item[主] 価値の信号と行動の対応から学習された自分の予測傾向が，他者側の予測に反映されること（転用）を示し，その成立条件を切り分けた．同じ軌跡で価値の信号と行動の対応を壊すと，自分側の快・不快への依存と他者側の符号に従う応答は0前後になり，他者側の符号応答も大きく弱まった．他者の状態を追う力が育つだけでは転用は生じず，学習後に快・不快の入口の重みだけを写しても戻らなかった．転用には，価値の信号と行動の一貫した対応と，自分の経験で学習された入口以外の重みや計算が要ると考えられる．
  \item[補助1] この転用は，他者の行動の予測を一般に正確にするものではない．本当の快・不快を与えたときの誤差は経験によって小さくならず，他者の行動の仕方が自分と違うと，経験の少ないモデルより外れることがあった．他者の規則の確率を自分の学習条件に近づけると，本当の快・不快の役立ち方が変わったが，位置の分布も同時に変わるので，自分の行動の仕方の投影という解釈は可能性にとどまる．
  \item[補助2] この転用は，他者固有の好みの自発的な表現や推定を保証しない．他者用のモジュールは好みではなく位置や方角を書き出し，好みの候補を入れ比べる照合も，他者の進んだ方角だけで選ぶ基準をほとんど上回らなかった．
\end{description}

なお，成立条件を切り分けた対照は，自己位置の追跡の事前基準を満たさなかった経験80\%条件で行った．主張の強さへの影響は\ref{sec:disc-strength}節で論じる．

以下，\ref{sec:related}章で関連研究を整理し，\ref{sec:method}章でモデルと実験方法を，\ref{sec:results}章で結果を述べる．
\ref{sec:discussion}章でシミュレーション説との関係と限界を論じ，\ref{sec:conclusion}章でまとめる．
